% Documentación científica editable de Arii.
% Compilar dos veces para actualizar el índice, referencias y enlaces.
\documentclass[11pt,a4paper,openany]{report}
\usepackage[utf8]{inputenc}
\usepackage[T1]{fontenc}
\usepackage[spanish,es-nodecimaldot]{babel}
\usepackage{lmodern,microtype,geometry,graphicx,xcolor}
\usepackage{amsmath,amssymb,mathtools,booktabs,tabularx}
\usepackage{enumitem,fancyhdr,titlesec,hyperref,bookmark,csquotes}

\geometry{top=24mm,bottom=23mm,inner=24mm,outer=21mm,headheight=15pt}
\definecolor{AriiNavy}{HTML}{071B3A}
\definecolor{AriiBlue}{HTML}{087DD1}
\definecolor{AriiCyan}{HTML}{11C6DF}
\definecolor{AriiViolet}{HTML}{7C3AED}
\definecolor{AriiMagenta}{HTML}{D946EF}
\definecolor{AriiGray}{HTML}{46566A}
\newcommand{\AriiVersion}{1.1.0}
\newcommand{\mixOmics}{\textit{mixOmics}}
\newcommand{\code}[1]{\texttt{#1}}
\newcommand{\RtwoX}{R^{2}_{X}}
\newcommand{\RtwoY}{R^{2}_{Y}}
\newcommand{\QtwoY}{Q^{2}_{Y}}
\newcommand{\vect}[1]{\boldsymbol{#1}}
\DeclareMathOperator*{\argmin}{arg\,min}
\IfFileExists{../assets/branding/readme-header-light.png}
  {\newcommand{\AriiHeader}{../assets/branding/readme-header-light.png}}
  {\IfFileExists{assets/branding/readme-header-light.png}
    {\newcommand{\AriiHeader}{assets/branding/readme-header-light.png}}
    {\newcommand{\AriiHeader}{readme-header-light.png}}}

\hypersetup{colorlinks=true,linkcolor=AriiBlue,urlcolor=AriiViolet,
  citecolor=AriiMagenta,pdftitle={Arii -- Documentación científica},
  pdfauthor={Gabriel Anderson}}
\setlist[itemize]{leftmargin=6mm,itemsep=2pt,topsep=4pt}
\setlist[enumerate]{leftmargin=7mm,itemsep=3pt,topsep=4pt}
\setcounter{tocdepth}{2}
\setcounter{secnumdepth}{2}
\titleformat{\chapter}[display]{\normalfont\bfseries\color{AriiNavy}}
  {\color{AriiBlue}\Large\chaptername\ \thechapter}{4pt}{\Huge}
\titlespacing*{\chapter}{0pt}{35pt}{24pt}
\titleformat{\section}{\normalfont\Large\bfseries\color{AriiNavy}}{\thesection}{0.7em}{}
\titleformat{\subsection}{\normalfont\large\bfseries\color{AriiNavy}}{\thesubsection}{0.7em}{}
\pagestyle{fancy}
\fancyhf{}
\fancyhead[L]{\small\textcolor{AriiGray}{Arii}}
\fancyhead[R]{\small\textcolor{AriiGray}{\nouppercase{\leftmark}}}
\fancyfoot[C]{\small\textcolor{AriiGray}{\thepage}}
\renewcommand{\headrulewidth}{0.3pt}
\renewcommand{\chaptermark}[1]{\markboth{#1}{}}
\fancypagestyle{plain}{
  \fancyhf{}
  \fancyfoot[C]{\small\textcolor{AriiGray}{\thepage}}
  \renewcommand{\headrulewidth}{0pt}
}
\raggedbottom

\begin{document}
\begin{titlepage}
  \centering\vspace*{13mm}
  \includegraphics[width=0.92\textwidth]{\AriiHeader}
  \vspace{22mm}
  {\Huge\bfseries\color{AriiNavy} Documentación científica de Arii\par}
  \vspace{8mm}
  {\Large\color{AriiGray}Fundamentos, implementación, validación y limitaciones\par}
  \vfill
  {\large Gabriel Anderson\par}\vspace{3mm}
  {\large Arii \AriiVersion\par}\vspace{15mm}
\end{titlepage}

\pagenumbering{roman}
\pdfbookmark[0]{Índice general}{toc}
\tableofcontents\clearpage\pagenumbering{arabic}

\chapter{Alcance y principio de diseño}

Arii implementa un flujo trazable para el análisis de una matriz ómica mediante
PCA, PLS-DA y una representación OPLS-DA binaria. La aplicación separa tres
niveles que no deben confundirse: el tratamiento técnico anterior a la matriz,
el preprocesamiento estadístico aprendido por el modelo y la presentación de los
resultados.

Los modelos base se calculan con \href{https://mixomics.org/}{mixOmics}. Arii añade el contrato de datos, la
validación agrupada, la selección de complejidad, diagnósticos, persistencia y una
transformación ortogonal propia del espacio latente para la vista OPLS-DA. Esta
última se documenta explícitamente porque no equivale a llamar a una función
OPLS-DA nativa de \mixOmics.

El documento describe qué calcula la versión \AriiVersion, qué supuestos exige y
qué conclusiones no permite. No reemplaza el criterio analítico específico de
RMN, GC/LC-MS, proteómica, transcriptómica o genómica.

\chapter{Contrato de datos y preprocesamiento}
\section{Matriz y metadatos}
El núcleo trabaja con
\begin{equation}
  \vect{X}\in\mathbb{R}^{n\times p},
\end{equation}
donde $n$ es el número de muestras y $p$ el número de variables. El eje ppm o
los identificadores son metadatos de variables; clase e individuo son
metadatos de muestras. El vector ppm nunca se incorpora como predictor.

El importador acepta las orientaciones \emph{variables por muestras} y
\emph{muestras por variables}, pero el motor siempre recibe $\vect{X}$ con la
forma anterior. La orientación se propone a partir de identificadores y
dimensiones; que una propuesta sea sintácticamente plausible no garantiza que
sea experimentalmente correcta. El usuario debe verificar nombres y
dimensiones antes de modelar.

Arii distingue un perfil continuo de una tabla de características. En el primer
caso existe un eje numérico y ordenado cuya vecindad tiene significado; en el
segundo, el orden de las variables no implica continuidad. Esta distinción
modifica la representación gráfica y la disponibilidad de STOCSY, pero no altera
la forma matricial del ajuste.

Las matrices deben llegar con el tratamiento técnico pertinente ya realizado:
normalización, alineamiento, corrección de deriva o lote, filtrado de calidad e
imputación cuando correspondan. Un cero analítico, una ausencia de detección y
un valor faltante no son equivalentes.

Cuando existe un identificador de individuo, éste define la unidad de
remuestreo: sus réplicas permanecen juntas en entrenamiento o validación.

El motor rechaza valores ausentes, infinitos o no numéricos. Arii no imputa de
forma automática porque la ausencia de detección, un cero analítico y un valor
faltante tienen significados diferentes según la plataforma. También elimina del
ajuste las variables con desviación estándar nula o numéricamente degenerada y
registra cuántas fueron excluidas.

\section{Centrado y escalado}
Sea $x_{ij}$ el valor de la variable $j$ en la muestra $i$, $\bar{x}_{j}$ su
media y $s_j$ su desviación estándar muestral. Arii ofrece:
\begin{align}
  x^{\mathrm{Pareto}}_{ij} &= \frac{x_{ij}-\bar{x}_{j}}{\sqrt{s_j}}, \\
  x^{\mathrm{UV}}_{ij} &= \frac{x_{ij}-\bar{x}_{j}}{s_j}, \\
  x^{\mathrm{centrado}}_{ij} &= x_{ij}-\bar{x}_{j}.
\end{align}
En validación supervisada, $\bar{x}_{j}$ y $s_j$ se estiman exclusivamente en
el entrenamiento y se aplican al conjunto externo para evitar fuga de
información.

Pareto reduce la influencia de las variables de gran escala sin igualarlas por
completo; el escalado de varianza unitaria concede un peso comparable a todas
las variables no constantes y puede amplificar ruido de pequeña varianza. La
opción centrada conserva las diferencias originales de escala. Ninguna elección
es universal: debe justificarse en relación con la plataforma y el objetivo.

\chapter{Análisis de componentes principales}
PCA aproxima la matriz centrada y escalada mediante
\begin{equation}
  \vect{X}=\vect{T}\vect{P}^{\mathsf T}+\vect{E},
\end{equation}
donde $\vect{T}$ contiene los scores, $\vect{P}$ los loadings y $\vect{E}$ los
residuos. Para $a$ componentes:
\begin{equation}
  \widehat{\vect{X}}_{a}=\vect{T}_{a}\vect{P}_{a}^{\mathsf T}.
\end{equation}
Arii informa
\begin{align}
  \mathrm{RMSE}_{X,a} &= \sqrt{\frac{1}{np}\sum_{i=1}^{n}\sum_{j=1}^{p}
    \left(x_{ij}-\widehat{x}_{ij,a}\right)^2}, \\
  \RtwoX(a) &= 1-\frac{\mathrm{SSE}_{X,a}}{\mathrm{TSS}_{X}}.
\end{align}
Son métricas descriptivas calculadas sobre las observaciones de ajuste; no son
rendimiento predictivo externo ni validan por sí solas el número de componentes.

Los scores son las coordenadas de las muestras en el espacio latente y los
loadings describen la dirección de cada componente en el espacio de variables.
La clase no interviene en el ajuste, aunque puede utilizarse para colorear la
figura. Por tanto, una agrupación visual por clase es un resultado descriptivo y
no una clasificación supervisada.

La cantidad máxima de componentes queda limitada por
$\min(n-1,p)$. Arii informa la fracción de variación de $\vect{X}$ explicada por
cada componente y su acumulado. Un aumento de $\RtwoX$ al añadir componentes es
esperable incluso si esas direcciones recogen ruido; por eso no existe un umbral
universal de varianza acumulada que valide el modelo.

\chapter{PLS-DA y OPLS-DA}
\section{PLS-DA}
PLS-DA relaciona $\vect{X}$ con una matriz indicadora
$\vect{Y}\in\{0,1\}^{n\times g}$ mediante variables latentes que maximizan su
covarianza. Arii ajusta \code{mixOmics::plsda} y asigna la clase mediante
\code{max.dist}. Las respuestas continuas se conservan para $\RtwoY$, $\QtwoY$, curvas de ROC y análisis de umbral.

La bondad de ajuste del modelo completo es
\begin{equation}
  \RtwoY = 1-\frac{\sum_i\lVert\vect{y}_i-\widehat{\vect{y}}_i\rVert^2}
  {\sum_i\lVert\vect{y}_i-\bar{\vect{y}}\rVert^2},
\end{equation}
y la capacidad predictiva interna utiliza predicciones externas:
\begin{equation}
  \QtwoY = 1-\frac{\mathrm{PRESS}_{Y}}{\mathrm{TSS}_{Y,\mathrm{CV}}}.
\end{equation}
Arii distingue siempre ajuste y CV. El primero aplica el modelo completo sobre
las mismas observaciones que lo estimaron y es descriptivo y potencialmente
optimista. El segundo agrega exclusivamente predicciones de muestras externas a
cada ajuste y aproxima rendimiento interno fuera de fold.

En $\QtwoY$, PRESS suma los errores externos. El denominador se centra respecto
de las proporciones de clase aprendidas en cada entrenamiento, evitando usar
información del conjunto externo para definir la referencia. Un valor alto no
garantiza transferencia a otro estudio: las particiones siguen procediendo del
mismo conjunto experimental.

\section{Clasificación y error balanceado}
\begin{align}
  \mathrm{Accuracy} &= \frac{1}{n}\sum_{i=1}^{n}
    \mathbb{I}(y_i=\widehat{y}_i), \\
  \mathrm{BER} &= \frac{1}{g}\sum_{k=1}^{g}\frac{\mathrm{FN}_{k}}{n_k}.
\end{align}
BER otorga el mismo peso a cada clase y es preferible cuando los tamaños de los
grupos difieren.

Arii agrega una matriz de confusión con predicciones externas y calcula AUC
one-vs-rest a partir de las respuestas continuas. La AUC macro es la media no
ponderada de las AUC por clase, por lo que cada grupo contribuye por igual. Las
métricas deben interpretarse conjuntamente: exactitud y AUC pueden ocultar una
clase con baja sensibilidad, mientras que BER resume el error medio por clase.

\newpage

\section{OPLS-DA}
Arii transforma un PLS-DA binario calculado con \mixOmics. Para un modelo de
$a$ variables latentes, toma sus scores originales $\vect{T}_a$ y la respuesta
continua predicha para la segunda clase. Tras centrar y normalizar esa respuesta,
obtiene una base predictiva unitaria $\vect{q}_p$. El residuo de los scores es

\begin{equation}
  \vect{T}_{\perp}=\vect{T}_a-
  \vect{q}_p\vect{q}_p^{\mathsf T}\vect{T}_a.
\end{equation}

Una descomposición QR de $\vect{T}_{\perp}$ produce hasta $a-1$ direcciones
ortogonales. Las bases se reescalan con la norma de los scores originales y los
loadings se reestiman por mínimos cuadrados:

\begin{equation}
  \vect{P}_{\mathrm{rot}}^{\mathsf T}=
  (\vect{T}_{\mathrm{rot}}^{\mathsf T}\vect{T}_{\mathrm{rot}})^{-1}
  \vect{T}_{\mathrm{rot}}^{\mathsf T}\vect{X}.
\end{equation}

La rotación no reajusta la respuesta y no cambia $\RtwoY$, $\QtwoY$, BER,
confusión, ROC ni las predicciones. Las componentes ortogonales describen
variación de $\vect{X}$ no asociada directamente con $\vect{Y}$ y no deben
interpretarse como biomarcadores discriminantes.

La rotación se calcula por separado para cada complejidad candidata. Por ello,
\emph{Ortogonal 1} de un modelo de dos componentes no debe asumirse idéntica a
\emph{Ortogonal 1} de un modelo de tres. Esta implementación es una vista
ortogonalizada y transparente del PLS-DA binario; no pretende reproducir de
forma bit a bit un algoritmo propietario ni una implementación OPLS externa.

\chapter{Validación cruzada y selección}
\section{Random subsets repetido}
La configuración predeterminada ejecuta $R=20$ repeticiones de $K=5$ folds
estratificados:
\begin{equation}
  N_{\mathrm{submodelos}}=R\times K=100.
\end{equation}
La fracción teórica por submodelo es 80/20, con variaciones por redondeo,
balance de clase o agrupación de individuos. Cada unidad recibe una predicción
externa por repetición.

La estratificación se realiza sobre unidades de remuestreo. Si existe un
identificador biológico, todas sus réplicas comparten fold; si no existe, cada
muestra es una unidad. Una clase debe contener al menos tantas unidades
independientes como folds. Arii registra la fracción externa efectiva, que puede
diferir del valor teórico cuando las unidades tienen distinto número de réplicas.

\section{Otros esquemas}
En \textbf{Monte Carlo}, cada iteración genera una partición aleatoria
estratificada independiente; Arii registra su cobertura efectiva.
\textbf{Leave-one-out} retira una unidad de remuestreo por submodelo, o el grupo
completo si existen individuos. \textbf{Venetian blinds} distribuye
sistemáticamente unidades ordenadas y sólo es defendible si el orden no codifica
clase, lote o deriva de forma favorable.

Estas estrategias estiman propiedades distintas y no son intercambiables por
conveniencia. Monte Carlo puede dejar algunas unidades menos representadas como
externas; leave-one-out utiliza conjuntos de entrenamiento muy semejantes y
puede presentar alta varianza; Venetian blinds depende de que el orden original
no introduzca información espuria.

Toda operación aprendida a partir de $\vect{X}$ debe ocurrir dentro del fold.
Arii aplica este principio al centrado, escalado y eliminación de variables
constantes. El filtrado supervisado, la imputación o la corrección de lote hecha
con el dataset completo antes de la validación pueden seguir produciendo fuga de
información aunque el motor de Arii procese correctamente sus folds.

\section{Selección de complejidad}
Sea $\overline{e}_{a}$ el BER CV medio y $\mathrm{SE}_{a}$ su error estándar. Si
$a_{\min}$ minimiza BER, Arii sugiere el menor $a$ que satisface
\begin{equation}
  \overline{e}_{a}\leq
  \overline{e}_{a_{\min}}+\mathrm{SE}_{a_{\min}}.
\end{equation}
Es una selección exploratoria. Una estimación final sin sesgo requiere
validación anidada o una cohorte externa.

Arii separa la \textbf{sugerencia} del \textbf{modelo actual}. Al terminar, la
sugerencia inicial se adopta como modelo actual, pero el usuario puede elegir
otra complejidad. Scores, loadings, confusión, ROC, exportación y permutaciones
se refieren al modelo actual; las métricas de los candidatos permanecen para
comparación. Esta distinción evita presentar una decisión manual como si fuera el
resultado automático de la regla.

\chapter{Resultados gráficos}
\section{Scores y T\texorpdfstring{$^2$}{2} de Hotelling}
Para una representación bidimensional con $q=2$ y $n$ observaciones:
\begin{equation}
  T^2_{\alpha}=\frac{q(n-1)}{n-q}F_{\alpha}(q,n-q).
\end{equation}
La elipse representa una región global de scores; no es un intervalo por clase
ni demuestra separación estadística.

La expresión se aplica a la proyección bidimensional visible y presupone una
estructura aproximadamente elíptica. Un punto fuera del límite es una señal para
revisión analítica, no evidencia automática de error ni autorización para excluir
la muestra.

\section{Loadings y gráfico de valores de carga}
El Loading plot aplica un reescalado dependiente del preprocesamiento:
\begin{align}
  v_{jk}^{\mathrm{Pareto}} &= \sqrt{s_j}\,p_{jk}, \\
  v_{jk}^{\mathrm{UV}} &= s_j\,p_{jk}, \\
  v_{jk}^{\mathrm{centrado}} &= p_{jk}.
\end{align}
El eje vertical representa $v_{jk}$ y el color el loading original con signo.
En RMN el eje ppm se dibuja descendente. El signo de una componente latente es
indeterminado y puede invertirse junto con sus scores sin alterar el modelo.

Para tablas de características Arii no conecta variables adyacentes con una
curva, porque esa continuidad sería artificial. Para perfiles numéricos sí
conserva la geometría del eje y, en RMN, lo dibuja en orden descendente.

\section{Curvas de ROC y umbral}
El AUC macro one-vs-rest es
\begin{equation}
  \mathrm{AUC}_{\mathrm{macro}}=\frac{1}{g}\sum_{k=1}^{g}\mathrm{AUC}_k.
\end{equation}
El umbral exploratorio maximiza el índice de Youden
\begin{equation}
  J(c)=\mathrm{Sensibilidad}(c)+\mathrm{Especificidad}(c)-1.
\end{equation}
No recalcula las métricas ni constituye un umbral clínicamente validado.

Las curvas se calculan por separado para el ajuste completo y para la media de
las respuestas externas disponibles por muestra. Esta separación evita que una
curva de ajuste optimista se presente como rendimiento validado.

\chapter{Tests de permutaciones}
Las etiquetas se permutan entre unidades intercambiables, manteniendo juntas las
réplicas de un individuo y los tamaños de clase. Los folds se reutilizan; como
$\vect{X}$ no cambia, también se reutilizan transformaciones permitidas.

La intercambiabilidad es un supuesto científico. Si existen bloques, pares,
familias, centros o series temporales, la permutación debe respetar esa
estructura. Informar un identificador de individuo protege réplicas del mismo
sujeto, pero no modela automáticamente todos los diseños experimentales posibles.

Para una estadística cuyo valor alto favorece el modelo:
\begin{equation}
  p_{\mathrm{emp}}=\frac{1+b}{B+1},
\end{equation}
donde $B$ es el número de permutaciones y $b$ los valores nulos al menos tan
extremos como el observado. La cola es superior para $\RtwoY$, $\QtwoY$ y AUC,
e inferior para BER. La resolución mínima es
\begin{equation}
  p_{\min}=\frac{1}{B+1}.
\end{equation}
La comparación residual y las gráficas frente a correlación de Y son
diagnósticos complementarios; no sustituyen el p empírico.

La complejidad permanece fija en la del modelo actual. Esto responde a la
pregunta: «¿un modelo con esta complejidad supera lo esperable bajo etiquetas
intercambiadas?». No corrige el optimismo de haber elegido esa complejidad con
los mismos datos. Una prueba que incluyera selección completa dentro de cada
permutación respondería una pregunta diferente y tendría un coste mayor.

En la comparación residual, la unidad del contraste es la muestra y su referencia
es el error medio de esa muestra bajo las permutaciones. Wilcoxon, la prueba de
signos y el test $t$ por aleatorización son diagnósticos relacionados; no deben
tratarse como tres oportunidades independientes de declarar significancia.

\chapter{STOCSY}
Sea $\vect{x}_{d}$ la variable pivote y $\vect{x}_{j}$ cualquier otra variable:
\begin{align}
  \operatorname{cov}(\vect{x}_{d},\vect{x}_{j}) &=
  \frac{1}{n-1}\sum_{i=1}^{n}(x_{id}-\bar{x}_{d})(x_{ij}-\bar{x}_{j}), \\
  r_{dj} &= \frac{\operatorname{cov}(\vect{x}_{d},\vect{x}_{j})}{s_d s_j}.
\end{align}
El perfil vertical representa covarianza y el color correlación. Arii usa todas
las muestras del dataset cargado, no sólo las del modelo visible. Una asociación
puede reflejar el mismo metabolito, covariación biológica, clase, lote o
confusión; STOCSY no demuestra causalidad ni identifica por sí solo un compuesto.

La selección gráfica se asigna a la variable cuyo eje es más cercano al punto
elegido. Si la variable pivote no presenta variación entre muestras, la
correlación no está definida y Arii rechaza el cálculo. El uso del dataset
completo aumenta el tamaño muestral, pero también significa que una selección de
muestras del modelo no restringe STOCSY.

\chapter{Ejecución y reproducibilidad}
El mismo worker R y especificación JSON se utilizan localmente y en Slurm. Las
validaciones y permutaciones independientes pueden distribuirse sin cambiar el
contrato estadístico. La GPU no participa actualmente y Arii no la consulta.

El worker, la especificación JSON y el resultado forman el límite reproducible
entre la interfaz y el motor. La ejecución paralela distribuye folds o
permutaciones independientes; las semillas se derivan de forma determinista
para que cambiar la cantidad de CPU no cambie la secuencia científica esperada.

Un proyecto \code{.arii} conserva configuración y resultados, pero referencia el
dataset mediante ruta y SHA-256 en lugar de incorporarlo. Esta decisión evita
duplicar matrices grandes o sensibles, a costa de exigir que el archivo original
permanezca disponible y sin modificaciones.

Un informe reproducible debe conservar dataset y huella, muestras, clases e
individuos, pretratamiento, centrado/escalado, modelo y complejidad, método CV,
semilla, unidad de remuestreo, versiones, métricas externas, permutaciones,
advertencias y validación independiente disponible.

\chapter{Limitaciones y uso responsable}

Los modelos con $p\gg n$ son especialmente vulnerables al sobreajuste. Ninguna
separación visual, $\RtwoY$ alto o prueba de permutación sustituye diseño
experimental, control de calidad, plausibilidad biológica y validación externa.

Arii no realiza actualmente imputación, corrección de lote, selección dispersa de
variables ni validación anidada. Si esos pasos se realizan antes de Arii usando
información de clase o del dataset completo, el rendimiento CV puede quedar
sesgado. El programa tampoco transforma automáticamente conteos,
composiciones o datos de célula única a una geometría apropiada.

La vista OPLS-DA es una rotación documentada del PLS-DA binario de Arii. Debe
citarse y describirse como tal; no corresponde atribuirle propiedades de una
implementación OPLS distinta. Del mismo modo, el umbral de Youden, la elipse de
Hotelling y STOCSY son ayudas exploratorias, no decisiones clínicas ni pruebas
de identificación molecular.

\begin{thebibliography}{9}
\bibitem{rohart2017} Rohart F, Gautier B, Singh A, Lê Cao K-A.
mixOmics: An R package for 'omics feature selection and multiple data integration.
\emph{PLoS Computational Biology}. 2017;13(11):e1005752.
\href{https://doi.org/10.1371/journal.pcbi.1005752}{doi:10.1371/journal.pcbi.1005752}.

\bibitem{wiklund2008} Wiklund S, Johansson E, Sjöström L, et al.
Visualization of GC/TOF-MS-based metabolomics data for identification of
biochemically interesting compounds using OPLS class models.
\emph{Analytical Chemistry}. 2008;80:115--122.
\href{https://doi.org/10.1021/ac0713510}{doi:10.1021/ac0713510}.

\bibitem{trygg2002} Trygg J, Wold S.
Orthogonal projections to latent structures (O-PLS).
\emph{Journal of Chemometrics}. 2002;16(3):119--128.
\href{https://doi.org/10.1002/cem.695}{doi:10.1002/cem.695}.

\bibitem{westerhuis2008} Westerhuis JA, Hoefsloot HCJ, Smit S, et al.
Assessment of PLSDA cross validation.
\emph{Metabolomics}. 2008;4(1):81--89.
\href{https://doi.org/10.1007/s11306-007-0099-6}{doi:10.1007/s11306-007-0099-6}.

\bibitem{cloarec2005} Cloarec O, Dumas M-E, Craig A, et al.
Statistical total correlation spectroscopy: an exploratory approach for latent
biomarker identification from metabolic $^1$H NMR data sets.
\emph{Analytical Chemistry}. 2005;77:1282--1289.
\href{https://doi.org/10.1021/ac048630x}{doi:10.1021/ac048630x}.
\end{thebibliography}
\end{document}
